\chapter{实验记录自动化：缺的不是工具，是默认行为}
\label{ch:logging}

\section{故事引入：三个月后的"考古工作"}

论文被接收了，但审稿人要求提供补充材料，解释表4中某个实验的具体设置。你打开代码仓库，开始"考古"：

\textbf{第一步：找日志}

你记得这个实验是三个月前跑的，打开 outputs/ 目录，看到一堆以日期命名的文件夹。但问题是，你记不清具体日期了。你只能一个一个打开，看里面的结果是不是那个实验。

\textbf{第二步：找配置}

终于找到了结果文件，但没有记录配置。你翻遍代码历史，试图找到当时用的超参数。你在某个 commit 里找到了可能的配置，但不确定是不是最终版本——你记得当时临时改过学习率，但不确定改成了多少。

\textbf{第三步：找数据}

代码里写的数据路径是 \texttt{data/v2/}，但你现在的数据目录是 \texttt{data/v3/}。你不记得当时是不是换过数据版本。你在聊天记录里搜索"数据"，试图找到线索。

\textbf{第四步：放弃}

折腾了一个下午，你决定重新跑一次这个实验。但由于参数不确定，重跑的结果和论文里报告的不一致。你只能在补充材料里写："由于时间久远，部分实验细节可能有偏差。"

\textbf{审稿人的回复：}"We cannot accept a paper where the authors cannot reproduce their own results."

\paragraph{这个悲剧本可以避免。}

如果你在实验结束时花 2 分钟记录关键信息，就不会有三个月后的噩梦。

问题不是缺少工具（MLflow、W\&B、TensorBoard 都很好），而是\textbf{缺少记录的默认行为}：很多人觉得"这次只是试试，不用记录"，结果试着试着就忘了记录，最后有价值的实验也没留下痕迹。

\section{两层日志策略：机器精确 + 人类简洁}

实验记录的核心挑战是平衡两个需求：
\begin{itemize}
  \item \textbf{机器需要完整精确的信息}（用于复现和自动化分析）；
  \item \textbf{人类需要简洁可读的摘要}（用于快速回顾和决策）。
\end{itemize}

如果只有机器日志（如 JSON），人类很难快速了解"这次实验到底想验证什么"；如果只有人类日志（如笔记），机器无法自动复现和对比。

\textbf{解决方案：两层日志，各司其职。}

\subsection{第一层：机器日志（run.json）}

\textbf{目的：}为复现和自动化提供完整、结构化的信息。

\textbf{原则：}
\begin{itemize}
  \item \textbf{自动生成：}不依赖人工输入，通过脚本自动收集；
  \item \textbf{结构化：}JSON 格式，便于程序解析和查询；
  \item \textbf{完整性：}包含复现所需的所有关键信息。
\end{itemize}

\textbf{最小字段集}（可直接照抄）：

\begin{verbatim}
{
  "run_id": "2026-02-01_1630_ablation_lr",
  "timestamp": {
    "start": "2026-02-01T16:30:45",
    "end": "2026-02-01T18:45:12"
  },
  "git": {
    "commit": "a1b2c3d4e5f6",
    "dirty": false,
    "branch": "exp/ablation-lr",
    "remote": "git@github.com:user/project.git"
  },
  "config": {
    "path": "configs/ablation_lr.yaml",
    "hash": "sha256:abcd1234...",
    "resolved": {
      "model": "transformer",
      "learning_rate": 3e-4,
      "batch_size": 32,
      ...
    }
  },
  "data": {
    "name": "dataset_v3",
    "path": "/data/project/v3",
    "hash": "sha256:ef567890...",
    "split": {
      "train": 8000,
      "val": 1000,
      "test": 1000
    }
  },
  "environment": {
    "python": "3.11.7",
    "cuda": "12.1",
    "platform": "Linux-5.15.0-x86_64",
    "gpu": "NVIDIA A100-SXM4-40GB",
    "pip_freeze_hash": "sha256:12345678..."
  },
  "random": {
    "seed": 42,
    "torch_seed": 42,
    "numpy_seed": 42,
    "python_seed": 42
  },
  "metrics": {
    "val_loss": 0.123,
    "val_acc": 0.943,
    "test_loss": 0.145,
    "test_acc": 0.931,
    "training_time_hours": 2.25
  },
  "artifacts": {
    "model": "outputs/2026-02-01_1630_ablation_lr/model.pt",
    "logs": "outputs/2026-02-01_1630_ablation_lr/train.log",
    "predictions": "outputs/2026-02-01_1630_ablation_lr/predictions.json",
    "plots": "outputs/2026-02-01_1630_ablation_lr/plots/"
  }
}
\end{verbatim}

\paragraph{关键字段解释：}

\begin{itemize}
  \item \textbf{run\_id：}唯一标识，建议用时间戳+简短描述（见第2章）。

  \item \textbf{git.commit：}代码版本，用 \texttt{git rev-parse HEAD} 获取。

  \item \textbf{git.dirty：}是否有未提交的修改，用 \texttt{git diff-index --quiet HEAD --} 检查。如果为 true，建议保存 diff：\texttt{git diff > changes.patch}。

  \item \textbf{config.resolved：}展开后的最终配置，包含所有默认值。这很重要，因为默认值可能随代码更新而变化。

  \item \textbf{data.hash：}数据版本的哈希值，确保数据完全一致。可以用 \texttt{sha256sum} 计算整个数据目录的哈希，或使用 DVC 等工具。

  \item \textbf{environment.pip\_freeze\_hash：}依赖版本的哈希，用 \texttt{pip freeze | sha256sum} 计算。避免完整保存 pip freeze 输出（太长），只保存哈希和原始文件路径。

  \item \textbf{random.seed：}所有随机数种子。确保设置了 PyTorch、NumPy、Python 内置 random 的种子。
\end{itemize}

\subsection{第二层：人类日志（run.md）}

\textbf{目的：}为人类（包括未来的你）提供快速理解实验的摘要。

\textbf{原则：}
\begin{itemize}
  \item \textbf{简洁：}不超过 10 行，核心信息一目了然；
  \item \textbf{结构化：}用固定的五要素组织；
  \item \textbf{手写：}允许主观判断和洞察。
\end{itemize}

\textbf{五要素模板}（5 行就够）：

\begin{verbatim}
# Run: 2026-02-01_1630_ablation_lr

## 假设 (Hypothesis)
测试学习率对收敛速度和最终性能的影响。预期：较小的学习率（1e-4）会更稳定。

## 变化 (Change)
相比 baseline（lr=3e-4），降低学习率至 1e-4，其他超参数不变。

## 结论 (Result)
- 收敛更慢（从 50 epoch 到 80 epoch）
- 最终性能略有提升（val_acc: 0.943 vs 0.938）
- 训练更稳定，loss 曲线无明显震荡

## 下一步 (Next)
尝试中间值 2e-4，可能平衡速度和性能。
考虑 learning rate warmup 策略。

## 风险/异常 (Risk)
无明显异常。数据增强可能需要配合调整（当前是固定的）。
\end{verbatim}

\paragraph{为什么只要 5 行？}

\begin{itemize}
  \item \textbf{降低记录门槛：}如果要写一篇长文档，你会拖延；5 行的话，2 分钟就能完成。
  \item \textbf{强制提炼核心：}逼迫你思考"这次实验到底验证了什么"，而不是流水账。
  \item \textbf{快速回顾：}几个月后再看，5 行摘要比完整日志更有用。
\end{itemize}

\section{自动化工具：让记录成为零成本行为}

\textbf{核心理念：}记录不应该依赖"记得去做"，而应该自动发生。

\subsection{在训练脚本中自动生成 run.json}

\paragraph{实现示例（Python）：}

\begin{verbatim}
import json
import subprocess
import hashlib
from pathlib import Path
from datetime import datetime

def log_run(run_id, config, metrics, output_dir):
    """
    自动记录实验信息到 run.json

    Args:
        run_id: 实验唯一标识
        config: 配置字典（已展开）
        metrics: 最终指标字典
        output_dir: 输出目录路径
    """
    run_info = {
        "run_id": run_id,
        "timestamp": {
            "start": datetime.now().isoformat(),
        },
        "git": get_git_info(),
        "config": {
            "resolved": config,
            "hash": hash_dict(config),
        },
        "data": get_data_info(config.get("data_path")),
        "environment": get_env_info(),
        "random": get_random_seeds(config),
        "metrics": metrics,
        "artifacts": {
            "model": str(output_dir / "model.pt"),
            "logs": str(output_dir / "train.log"),
        }
    }

    # 保存到文件
    run_file = output_dir / "run.json"
    with open(run_file, "w") as f:
        json.dump(run_info, f, indent=2)

    print(f"Run info logged to {run_file}")

def get_git_info():
    """获取 git 信息"""
    try:
        commit = subprocess.check_output(
            ["git", "rev-parse", "HEAD"]
        ).decode().strip()

        # 检查是否有未提交的修改
        subprocess.check_call(
            ["git", "diff-index", "--quiet", "HEAD", "--"]
        )
        dirty = False
    except subprocess.CalledProcessError:
        dirty = True

    branch = subprocess.check_output(
        ["git", "rev-parse", "--abbrev-ref", "HEAD"]
    ).decode().strip()

    remote = subprocess.check_output(
        ["git", "config", "--get", "remote.origin.url"]
    ).decode().strip()

    return {
        "commit": commit,
        "dirty": dirty,
        "branch": branch,
        "remote": remote
    }

def get_data_info(data_path):
    """获取数据信息"""
    data_path = Path(data_path)

    # 计算数据目录哈希（简化版，实际可用 DVC）
    # 这里只是示例，真实场景建议用专门的数据版本管理工具

    return {
        "name": data_path.name,
        "path": str(data_path.absolute()),
        # "hash": compute_dir_hash(data_path),  # 可选
    }

def get_env_info():
    """获取环境信息"""
    import sys
    import platform

    env = {
        "python": sys.version.split()[0],
        "platform": platform.platform(),
    }

    # 获取 CUDA 版本（如果可用）
    try:
        import torch
        env["cuda"] = torch.version.cuda
        env["pytorch"] = torch.__version__
        if torch.cuda.is_available():
            env["gpu"] = torch.cuda.get_device_name(0)
    except ImportError:
        pass

    # 保存 pip freeze 到单独文件
    pip_freeze = subprocess.check_output(
        ["pip", "freeze"]
    ).decode()
    pip_file = Path("requirements_freeze.txt")
    pip_file.write_text(pip_freeze)

    env["pip_freeze_hash"] = hashlib.sha256(
        pip_freeze.encode()
    ).hexdigest()[:16]

    return env

def get_random_seeds(config):
    """提取随机种子"""
    return {
        "seed": config.get("seed", None),
        "torch_seed": config.get("torch_seed", None),
        "numpy_seed": config.get("numpy_seed", None),
    }

def hash_dict(d):
    """计算字典的哈希"""
    import json
    return hashlib.sha256(
        json.dumps(d, sort_keys=True).encode()
    ).hexdigest()[:16]
\end{verbatim}

\paragraph{在训练脚本中使用：}

\begin{verbatim}
# train.py

import argparse
from pathlib import Path
from run_logger import log_run  # 上面的工具

def main():
    args = parse_args()

    # 创建 run_id 和输出目录
    run_id = f"{datetime.now().strftime('%Y-%m-%d_%H%M')}_{args.exp_name}"
    output_dir = Path("outputs") / run_id
    output_dir.mkdir(parents=True, exist_ok=True)

    # 加载和展开配置
    config = load_config(args.config)
    config = resolve_config(config, args)  # 展开所有默认值

    # 设置随机种子
    set_random_seeds(config["seed"])

    # 训练模型
    model, metrics = train_model(config, output_dir)

    # 自动记录实验信息
    log_run(
        run_id=run_id,
        config=config,
        metrics=metrics,
        output_dir=output_dir
    )

    # 提示写 run.md
    print(f"\n{'='*60}")
    print(f"[OK] Experiment completed: {run_id}")
    print(f"[NOTE] Please write a brief summary in:")
    print(f"    {output_dir / 'run.md'}")
    print(f"{'='*60}\n")

if __name__ == "__main__":
    main()
\end{verbatim}

\subsection{用模板简化 run.md 编写}

在 output 目录自动生成 \texttt{run.md} 模板：

\begin{verbatim}
def create_run_md_template(output_dir, run_id):
    """创建 run.md 模板"""
    template = f"""# Run: {run_id}

## 假设 (Hypothesis)
[这次实验想验证什么？预期结果是什么？]

## 变化 (Change)
[相比之前的实验，改了什么？]

## 结论 (Result)
[实验结果如何？有什么意外发现？]

## 下一步 (Next)
[基于这次结果，下一步打算做什么？]

## 风险/异常 (Risk)
[有没有什么值得注意的异常或风险？]
"""

    md_file = output_dir / "run.md"
    if not md_file.exists():
        md_file.write_text(template)
        print(f"[NOTE] run.md template created at {md_file}")
\end{verbatim}

这样，每次实验结束后，你只需要填空，而不是从零开始写。

\section{与现有工具集成}

\subsection{与 MLflow 集成}

如果你已经在使用 MLflow，可以把 run.json 信息同步到 MLflow：

\begin{verbatim}
import mlflow

def log_to_mlflow(run_info):
    """将 run.json 信息记录到 MLflow"""
    with mlflow.start_run(run_name=run_info["run_id"]):
        # 记录参数
        mlflow.log_params(run_info["config"]["resolved"])

        # 记录指标
        mlflow.log_metrics(run_info["metrics"])

        # 记录环境信息
        mlflow.log_dict(run_info["environment"], "environment.json")

        # 记录 git 信息
        mlflow.set_tag("git.commit", run_info["git"]["commit"])
        mlflow.set_tag("git.branch", run_info["git"]["branch"])
        mlflow.set_tag("git.dirty", run_info["git"]["dirty"])

        # 记录产物
        mlflow.log_artifact(run_info["artifacts"]["model"])
\end{verbatim}

\subsection{与 Weights \& Biases 集成}

\begin{verbatim}
import wandb

def log_to_wandb(run_info):
    """将 run.json 信息记录到 W&B"""
    wandb.init(
        project="my-research",
        name=run_info["run_id"],
        config=run_info["config"]["resolved"],
        tags=[run_info["git"]["branch"]]
    )

    # 记录指标
    wandb.log(run_info["metrics"])

    # 记录环境和 git 信息
    wandb.config.update({
        "git_commit": run_info["git"]["commit"],
        "git_dirty": run_info["git"]["dirty"],
        "python_version": run_info["environment"]["python"],
    })

    # 保存模型
    wandb.save(run_info["artifacts"]["model"])
\end{verbatim}

\textbf{重点：}工具是辅助，核心是\textbf{记录的规范}。即使没有 MLflow/W\&B，有 run.json 和 run.md 也足够。

\section{日志查询与分析}

有了结构化的 run.json，你可以快速查询和对比实验：

\subsection{查找最好的实验}

\begin{verbatim}
# find_best_run.py

import json
from pathlib import Path

def find_best_runs(metric="test_acc", top_k=5):
    """查找指标最好的实验"""
    runs = []

    for run_dir in Path("outputs").iterdir():
        if not run_dir.is_dir():
            continue

        run_json = run_dir / "run.json"
        if not run_json.exists():
            continue

        with open(run_json) as f:
            run_info = json.load(f)

        if metric in run_info.get("metrics", {}):
            runs.append({
                "run_id": run_info["run_id"],
                metric: run_info["metrics"][metric],
                "config": run_info["config"]["resolved"]
            })

    # 排序
    runs.sort(key=lambda x: x[metric], reverse=True)

    print(f"Top {top_k} runs by {metric}:")
    for i, run in enumerate(runs[:top_k], 1):
        print(f"{i}. {run['run_id']}: {run[metric]:.4f}")
        print(f"   Config: lr={run['config'].get('learning_rate')}, "
              f"bs={run['config'].get('batch_size')}")

    return runs[:top_k]

if __name__ == "__main__":
    find_best_runs()
\end{verbatim}

\subsection{对比两个实验的配置差异}

\begin{verbatim}
# compare_runs.py

import json
from pathlib import Path

def compare_runs(run_id1, run_id2):
    """对比两个实验的配置和结果"""
    run1 = load_run(run_id1)
    run2 = load_run(run_id2)

    print(f"Comparing {run_id1} vs {run_id2}\n")

    # 对比配置
    config1 = run1["config"]["resolved"]
    config2 = run2["config"]["resolved"]

    print("Configuration differences:")
    for key in set(config1.keys()) | set(config2.keys()):
        val1 = config1.get(key, "N/A")
        val2 = config2.get(key, "N/A")
        if val1 != val2:
            print(f"  {key}: {val1} -> {val2}")

    # 对比指标
    print("\nMetrics:")
    metrics1 = run1.get("metrics", {})
    metrics2 = run2.get("metrics", {})
    for key in set(metrics1.keys()) & set(metrics2.keys()):
        val1 = metrics1[key]
        val2 = metrics2[key]
        diff = val2 - val1
        print(f"  {key}: {val1:.4f} -> {val2:.4f} "
              f"({diff:+.4f}, {diff/val1*100:+.2f}%)")

def load_run(run_id):
    """加载实验信息"""
    run_json = Path("outputs") / run_id / "run.json"
    with open(run_json) as f:
        return json.load(f)

if __name__ == "__main__":
    import sys
    if len(sys.argv) != 3:
        print("Usage: python compare_runs.py <run_id1> <run_id2>")
        sys.exit(1)

    compare_runs(sys.argv[1], sys.argv[2])
\end{verbatim}

\section{常见问题与解决方案}

\subsection{Q1：run.json 太大了怎么办？}

\textbf{问题：}如果保存完整的 config（包含模型定义、数据预处理细节等），run.json 可能很大。

\textbf{解决方案：}
\begin{itemize}
  \item 只保存"关键超参数"在 run.json 里（如学习率、batch size）；
  \item 完整 config 保存到单独文件：\texttt{config\_resolved.yaml}；
  \item run.json 里只记录 config 文件路径和哈希。
\end{itemize}

\subsection{Q2：忘了记录 run.md 怎么办？}

\textbf{解决方案：}
\begin{itemize}
  \item 在 Makefile 或脚本里加入检查：
  \begin{verbatim}
  if [ ! -f outputs/$RUN_ID/run.md ]; then
      echo "Warning: run.md not found for $RUN_ID"
      echo "Please write a summary before continuing."
  fi
  \end{verbatim}

  \item 定期（如每周五）检查哪些实验缺少 run.md，集中补充。

  \item 实在想不起来就写"遗忘"，总比没有记录强。
\end{itemize}

\subsection{Q3：数据版本太大，无法计算哈希怎么办？}

\textbf{解决方案：}
\begin{itemize}
  \item 使用数据版本管理工具（如 DVC、Git LFS）；
  \item 或者只记录数据的"清单文件"（manifest）：
  \begin{verbatim}
  # 生成清单
  find data/ -type f | xargs sha256sum > data_manifest.txt

  # 计算清单哈希
  sha256sum data_manifest.txt
  \end{verbatim}

  \item 在 run.json 里记录清单文件路径和哈希。
\end{itemize}

\section{10 分钟动作：为下一个实验建立自动记录}

如果你现在只做一件事：建立自动记录的最小系统。

\begin{enumerate}
  \item \textbf{复制 run\_logger.py}

  把前面的 \texttt{log\_run} 函数复制到你的项目里。

  \item \textbf{在训练脚本末尾调用}
  \begin{verbatim}
  # 训练结束后
  log_run(run_id, config, metrics, output_dir)
  \end{verbatim}

  \item \textbf{创建 run.md 模板}
  \begin{verbatim}
  create_run_md_template(output_dir, run_id)
  \end{verbatim}

  \item \textbf{跑一次实验测试}

  跑一次实验，确认：
  \begin{itemize}
    \item outputs/<run\_id>/run.json 自动生成了
    \item outputs/<run\_id>/run.md 模板生成了
    \item run.json 里的信息完整（git、config、env、metrics）
  \end{itemize}

  \item \textbf{填写 run.md}

  花 2 分钟填写五要素模板。
\end{enumerate}

从下一个实验开始，记录就是自动的、零成本的。你唯一要做的就是花 2 分钟写 5 行摘要——这个投入会在三个月后得到百倍回报。
